\subsection{Introducción Teórica}
\label{subsec:introduccion-teorica}

% ####################################################################################################################
% ####################################################################################################################

\subsubsection{Mediciones e Incertezas}

Toda medición experimental contiene una incerteza inherente \cite{catedra2007mediciones}.
Esta incerteza puede ser de naturaleza \textbf{sistemática} (siempre del mismo signo, debida a calibración del instrumento, método inadecuado, etc.) o \textbf{accidental/casual} (de signo variable, debida a errores de apreciación del observador, fluctuaciones del objeto, etc).

En el caso de instrumentos digitales, la incerteza sistemática queda determinada por las especificaciones del fabricante, y se calcula a partir de dos contribuciones: el error proporcional a la lectura (\textit{reading}) y el error en dígitos (dgt) sobre la última cifra del display.

La expresión final de una medición directa es: $X = X_0 \pm \Delta X$.

% ####################################################################################################################
% ####################################################################################################################

\subsubsection{Ley de Ohm}

\todo{P1: Ley de Ohm ($V = IR$), relación lineal $V$ vs.\ $I$, ajuste $V = mI + b$ con $m = R$. Significado de $m$ y $b$; diferencia entre regresión y promedio de $R_n = V_n/I_n$.}

% ####################################################################################################################
% ####################################################################################################################

\subsubsection{Multímetro Digital como Ohmímetro}

El multímetro digital (tester) permite medir resistencias operando como \textbf{ohmímetro}: una fuente de corriente interna hace circular una corriente conocida $I$ a través de la resistencia desconocida $R$, y un voltímetro interno mide la caída de tensión $V$ resultante. Dado que $V = IR$, el valor de $R$ se obtiene como $R = V/I$.

El voltímetro interno tiene una lectura a fondo de escala de \textbf{200 mV}. La corriente de prueba varía según el rango seleccionado:

\begin{itemize}
    \item Rango 200 $\Omega$ $\rightarrow$ $I = 1$ mA \quad (1 mA $\times$ 200 $\Omega$ = 200 mV)
    \item Rango 2 k$\Omega$ $\rightarrow$ $I = 0{,}1$ mA
    \item Rangos mayores $\rightarrow$ corrientes menores
\end{itemize}

La incerteza de la medición se estima con la fórmula del fabricante:

\begin{equation}
    \Delta R = \frac{\%rdg}{100} \times R_{medida} + dgt \times resolución
    \label{eq:incerteza-multimetro}
\end{equation}

donde \textbf{\%rdg} es el error porcentual sobre la lectura (\textit{reading}) y \textbf{dgt} es la cantidad de dígitos de incerteza en la última cifra significativa.

% ####################################################################################################################
% ####################################################################################################################

\subsubsection{Código de Colores de Resistencias}

Las resistencias tienen 4 franjas de color: las dos primeras son dígitos, la tercera es el multiplicador y la cuarta es la tolerancia (dorado = $\pm 5\%$, plateado = $\pm 10\%$).

\begin{equation}
    R_{nominal} = (\text{dígito}_1 \cdot 10 + \text{dígito}_2) \times \text{multiplicador}
    \label{eq:codigo-colores}
\end{equation}

% ####################################################################################################################
% ####################################################################################################################

\subsubsection{Influencia del Aparato de Medida}

\todo{P3: 2ª Ley de Kirchhoff ($\sum V = 0$), divisor resistivo ideal $V_{le\acute{i}da} = \frac{R_2}{R_1+R_2}V_0$. Con voltímetro real ($R_{int} = 10\ \text{M}\Omega$): $R'_2 = R_2 \| R_{int}$; el efecto es apreciable cuando $R_2 \sim R_{int}$. ($R_1$, $R_2$ son las del circuito, no las de la Práctica 2.)}
